\documentclass[12pt,a4paper]{article}
\newcommand{\LabNumber}{3}
\newcommand{\LabTitle}{Cifruri polialfabetice:\\ Cifrul Vigenère pentru limba română}
\newcommand{\LabTitleShort}{Lucrarea de laborator nr. 3 — Criptografie \& Securitate Informațională}
% ---------------------------------------------------------------------------
% Common preamble for the CS laboratory reports (compile with XeLaTeX)
% Strunga Daniel-Ioan, FAF-242 — https://github.com/don09dandidi/CS
% ---------------------------------------------------------------------------
\usepackage{fontspec}
\setmainfont{Latin Modern Roman}
\setsansfont{Latin Modern Sans}
\setmonofont{DejaVu Sans Mono}[Scale=0.85]
\usepackage[english]{babel}
\babelprovide[import, main]{romanian}
\usepackage[a4paper, margin=2.5cm, headheight=15pt]{geometry}
\usepackage{amsmath, amssymb}
\usepackage{booktabs, array, longtable, tabularx}
\usepackage{xcolor}
\usepackage{listings}
\usepackage{fvextra}
\usepackage[most]{tcolorbox}
\usepackage{fancyhdr}
\usepackage{tikz}
\usetikzlibrary{arrows.meta, positioning, shapes.geometric, fit, calc}
\usepackage{enumitem}
\usepackage{caption}
\usepackage{float}
\usepackage[hidelinks]{hyperref}
\usepackage{url}

\definecolor{codebg}{RGB}{248,248,250}
\definecolor{codeframe}{RGB}{200,200,210}
\definecolor{kw}{RGB}{0,0,160}
\definecolor{cmt}{RGB}{0,120,0}
\definecolor{str}{RGB}{160,30,30}
\definecolor{termbg}{RGB}{30,32,40}
\definecolor{termfg}{RGB}{225,228,235}
\definecolor{termbar}{RGB}{60,63,75}
\definecolor{boxblue}{RGB}{230,238,250}

\lstdefinestyle{py}{
  language=Python, basicstyle=\ttfamily\small, backgroundcolor=\color{codebg},
  frame=single, rulecolor=\color{codeframe}, numbers=left, numberstyle=\tiny\color{gray},
  keywordstyle=\color{kw}\bfseries, commentstyle=\color{cmt}\itshape, stringstyle=\color{str},
  breaklines=true, showstringspaces=false, columns=fullflexible, keepspaces=true, tabsize=4,
  extendedchars=true, inputencoding=utf8
}
\lstdefinestyle{sh}{
  language=bash, basicstyle=\ttfamily\small, backgroundcolor=\color{codebg},
  frame=single, rulecolor=\color{codeframe}, numbers=left, numberstyle=\tiny\color{gray},
  keywordstyle=\color{kw}\bfseries, commentstyle=\color{cmt}\itshape, stringstyle=\color{str},
  breaklines=true, showstringspaces=false, columns=fullflexible, keepspaces=true
}
\lstdefinestyle{wl}{
  basicstyle=\ttfamily\small, backgroundcolor=\color{codebg},
  frame=single, rulecolor=\color{codeframe}, numbers=left, numberstyle=\tiny\color{gray},
  morekeywords={RandomPrime,RandomInteger,PowerMod,Mod,GCD,FromDigits,IntegerDigits,
    ToCharacterCode,FromCharacterCode,IntegerLength,BitLength,Module,Print,Hash,Length},
  keywordstyle=\color{kw}\bfseries, commentstyle=\color{cmt}\itshape,
  morecomment=[s]{(*}{*)}, breaklines=true, columns=fullflexible, keepspaces=true
}

% Console session box (as in the Lab 2 report)
\newtcolorbox{terminal}[1][Sesiune Terminal Consolă]{
  enhanced, breakable, colback=termbg, colframe=termbar, coltext=termfg,
  fonttitle=\small\ttfamily\bfseries, title={\textbullet\textbullet\textbullet\quad #1},
  left=4pt, right=4pt, top=2pt, bottom=2pt, boxrule=0.6pt, arc=2pt
}
\DefineVerbatimEnvironment{termout}{Verbatim}{fontsize=\scriptsize, breaklines=true, breakanywhere=true}

\newtcolorbox{infobox}[1]{
  enhanced, breakable, colback=boxblue, colframe=kw!60, fonttitle=\bfseries,
  title={#1}, boxrule=0.6pt, arc=2pt
}

\newcommand{\bignum}[1]{\par\nobreak\vspace{1pt}{\noindent\ttfamily\footnotesize\seqsplit{#1}\par}\vspace{5pt}}
\usepackage{seqsplit}

\pagestyle{fancy}
\fancyhf{}
\renewcommand{\headrulewidth}{0.4pt}
\fancyhead[L]{\footnotesize\textit{\LabTitleShort}}
\fancyhead[R]{\footnotesize strunga.daniel-ioan}
\fancyfoot[C]{\thepage}

\newcommand{\makereporttitle}{%
\begin{titlepage}
  \centering
  {\large\bfseries UNIVERSITATEA TEHNICĂ A MOLDOVEI\par}
  \vspace{2pt}
  {\normalsize Facultatea Calculatoare, Informatică și Microelectronică\par}
  {\normalsize Departamentul Ingineria Software și Automatică\par}
  \vspace{5cm}
  {\LARGE\bfseries LUCRAREA DE LABORATOR NR. \LabNumber\par}
  \vspace{8pt}
  {\Large\bfseries \LabTitle\par}
  \vspace{10pt}
  {\normalsize Disciplina: Criptografie și Securitate Informațională\par}
  \vfill
  \begin{tabular*}{\textwidth}{@{}l@{\extracolsep{\fill}}r@{}}
    \textbf{Autor:} & \textbf{Verificat de:} \\
    Daniel-Ioan Strunga & Lector / Supervizor \\
    Student gr. FAF-242 & Maia Zaica \\
  \end{tabular*}
  \vspace{3cm}

  {\normalsize Chișinău — 2026\par}
\end{titlepage}
}


\begin{document}
\makereporttitle
\tableofcontents
\newpage

\section{Introduction and Laboratory Objectives}
The weakness of every monoalphabetic cipher (Laboratory Work No.~2) is that the ciphertext
frequency distribution is a simple relabelling of the plaintext distribution. A cipher becomes
cryptographically stronger when its distribution is \emph{flattened}, i.e.\ when the same
plaintext letter can be encrypted into different ciphertext letters depending on its position.
Polyalphabetic ciphers achieve this by alternating between several substitution alphabets.

The most famous representative is the \textbf{Vigenère cipher}, first described by Giovan
Battista Bellaso in 1553 (\emph{La cifra del Sig.\ Giovan Battista Bellaso}) and later
misattributed to Blaise de Vigenère. Instead of a single Caesar shift it applies a sequence of
shifts $k_0,k_1,\dots,k_{L-1}$ given by the letters of a key word.

According to the group register, the student number is $n = 22$ (even), therefore the assigned
task is \textbf{Task 3.2}:
\begin{enumerate}
  \item Implement the Vigenère algorithm for messages in the \textbf{Romanian language} with an
        alphabet of \textbf{31 letters}, encoded as $0,1,\dots,30$.
  \item Accept only letters (\texttt{A}–\texttt{Z}, \texttt{a}–\texttt{z} and the Romanian
        diacritics); on any other symbol suggest the correct character range to the user.
  \item Enforce a minimum key length of \textbf{7} letters.
  \item Before processing, remove all spaces and convert all letters to upper case.
  \item Let the user choose the operation (encryption/decryption), enter the key and the
        message/cryptogram, and display the result, using exactly the formulas of the
        mathematical model.
\end{enumerate}

\section{Theoretical Foundations}
\subsection{Encoding of the Romanian Alphabet}
Let $\Sigma_{RO}$ be the Romanian alphabet with $n = 31$ symbols and let
$\mathbb{Z}_{31}=\{0,\dots,30\}$. The encoding used throughout the laboratory is given in
Table~\ref{tab:alpha}.

\begin{table}[H]
\centering
\caption{Encoding of the Romanian alphabet ($n = 31$)}
\label{tab:alpha}
\small
\setlength{\tabcolsep}{3.2pt}
\begin{tabular}{l*{16}{c}}
\toprule
Letter & A & Ă & Â & B & C & D & E & F & G & H & I & Î & J & K & L & M \\
Code   & 0 & 1 & 2 & 3 & 4 & 5 & 6 & 7 & 8 & 9 & 10 & 11 & 12 & 13 & 14 & 15 \\
\midrule
Letter & N & O & P & Q & R & S & Ș & T & Ț & U & V & W & X & Y & Z & \\
Code   & 16 & 17 & 18 & 19 & 20 & 21 & 22 & 23 & 24 & 25 & 26 & 27 & 28 & 29 & 30 & \\
\bottomrule
\end{tabular}
\end{table}

\subsection{Mathematical Model of the Vigenère Cipher}
Let the plaintext be $m=(m_0,\dots,m_{N-1})$ and the key $k=(k_0,\dots,k_{L-1})$ with
$m_i,k_j\in\mathbb{Z}_{31}$ and $L\ge 7$. The key is repeated periodically over the message:
\begin{align}
  c_i &= (m_i + k_{i \bmod L}) \bmod 31, \label{eq:enc}\\
  m_i &= (c_i - k_{i \bmod L}) \bmod 31. \label{eq:dec}
\end{align}
Decryption is the inverse operation because $(m_i + k - k) \bmod 31 = m_i$ for every
$k\in\mathbb{Z}_{31}$. Each key letter selects one of the 31 rows of the \emph{Tabula Recta},
so the cipher is a periodic sequence of $L$ Caesar ciphers.

\subsection{Key Space and Security}
For a key of length $L$ over 31 letters, the number of possible keys is $31^{L}$. With the minimum
admissible length $L = 7$:
\[
  |K| = 31^{7} = 27\,512\,614\,111 \approx 2.75\times10^{10},
\]
and for $L = 8$ already $|K|\approx 8.53\times10^{11}$. In addition, an attacker does not know
$L$ in advance. However, once $L$ is found (Kasiski examination or the Index of Coincidence of
every $L$-th letter, see Lab~2), the cryptogram splits into $L$ independent Caesar ciphers that
are broken separately by frequency analysis. Long keys (ideally as long as the message, i.e.\ a
running key / one-time pad) make the cipher considerably harder to attack.

\section{Architecture of the Python Application}
The application (\texttt{lab3/vigenere\_ro.py}) is written in Python~3 without external
dependencies. The alphabet is indexed through an explicit lookup table; no ASCII/Unicode
arithmetic (\texttt{ord}/\texttt{chr}) is used, so the diacritics are handled exactly like the
other letters.

\begin{infobox}{Application Component Structure}
\begin{itemize}[leftmargin=*]
  \item \textbf{Data layer:} \texttt{ROMANIAN\_ALPHABET} (31 letters) and the dictionary
        \texttt{CHAR\_TO\_CODE} implementing Table~\ref{tab:alpha}.
  \item \textbf{Input layer:} \texttt{normalize()} — unifies Ş/Ţ (cedilla) to Ș/Ț (comma below),
        removes all whitespace and converts to upper case; \texttt{validate()} and
        \texttt{prepare\_key()} reject invalid symbols and keys shorter than 7 letters, printing
        the allowed character range.
  \item \textbf{Cipher core:} \texttt{vigenere(text, key, decrypt)} implements
        formulas~\eqref{eq:enc}–\eqref{eq:dec}; with \texttt{trace=True} it prints the step
        table in the same format as the laboratory guide.
  \item \textbf{UI layer:} an interactive menu (encrypt / decrypt / exit) and a
        \texttt{--demo} mode with the deterministic examples used in this report.
\end{itemize}
\end{infobox}

\subsection{Cipher Core}
\begin{lstlisting}[style=py, caption={Core of the Vigenère implementation (excerpt)}]
def vigenere(text: str, key: str, decrypt: bool = False, trace: bool = False) -> str:
    sign = -1 if decrypt else 1
    out = []
    for i, ch in enumerate(text):
        m = CHAR_TO_CODE[ch]                    # code of the text letter
        k = CHAR_TO_CODE[key[i % len(key)]]     # key repeated periodically
        c = (m + sign * k) % N                  # N = 31
        out.append(ROMANIAN_ALPHABET[c])
    return "".join(out)
\end{lstlisting}

\subsection{Input Normalisation and Validation}
\begin{lstlisting}[style=py, caption={Normalisation and validation of the input}]
def normalize(text: str) -> str:
    text = (text.replace('Ş', 'Ș').replace('ş', 'ș')
                .replace('Ţ', 'Ț').replace('ţ', 'ț'))
    text = "".join(text.split())   # remove spaces
    return text.upper()

def prepare_key(key: str) -> str:
    key = normalize(key)
    validate(key, "cheia")         # only the 31 letters are allowed
    if len(key) < MIN_KEY_LEN:     # MIN_KEY_LEN = 7
        raise InvalidInputError(...)
    return key
\end{lstlisting}

\section{Experimental Results}
\subsection{Example 1: \textit{Per aspera ad astra}}
The example from the laboratory guide is repeated over the 31-letter alphabet. The key from the
guide (\texttt{SUPER}) is shorter than 7 letters and is rejected by the program, so the key
\texttt{SUPERMAN} ($L=8$) was used.

\begin{table}[H]
\centering
\caption{Encryption of \texttt{PERASPERAADASTRA} with $K=$ \texttt{SUPERMAN} (mod 31)}
\small
\setlength{\tabcolsep}{3.5pt}
\begin{tabular}{l*{16}{c}}
\toprule
Plaintext $M$ & P & E & R & A & S & P & E & R & A & A & D & A & S & T & R & A \\
Key $K$       & S & U & P & E & R & M & A & N & S & U & P & E & R & M & A & N \\
\midrule
$m_i$         & 18 & 6 & 20 & 0 & 21 & 18 & 6 & 20 & 0 & 0 & 5 & 0 & 21 & 23 & 20 & 0 \\
$k_i$         & 21 & 25 & 18 & 6 & 20 & 15 & 0 & 16 & 21 & 25 & 18 & 6 & 20 & 15 & 0 & 16 \\
$m_i+k_i \bmod 31$ & 8 & 0 & 7 & 6 & 10 & 2 & 6 & 5 & 21 & 25 & 23 & 6 & 10 & 7 & 20 & 16 \\
\midrule
Ciphertext $C$ & G & A & F & E & I & Â & E & D & S & U & T & E & I & F & R & N \\
\bottomrule
\end{tabular}
\end{table}

Result: $C=$ \texttt{GAFEIÂEDSUTEIFRN}. Note that the five occurrences of \texttt{A} are
encrypted as \texttt{E}, \texttt{S}, \texttt{U}, \texttt{E} and \texttt{N}, while the two
\texttt{E} letters produce \texttt{A} and \texttt{E} — the frequency profile of the plaintext is
no longer preserved. Applying formula~\eqref{eq:dec} with the same key returns
\texttt{PERASPERAADASTRA}.

\subsection{Example 2: Message with Romanian Diacritics}
Message: \emph{„Criptografia și securitatea informațională la UTM”}, key \texttt{CHIȘINĂU}
($L = 8$, contains diacritics):

\begin{terminal}
\begin{termout}
$ python3 vigenere_ro.py --demo
>>> Criptare: M = 'Criptografia și securitatea informațională la UTM'
    ->  'CRIPTOGRAFIAȘISECURITATEAINFORMAȚIONALĂLAUTM',  K = 'CHIȘINĂU'
  Textul clar M     C   R   I   P   T   O   G   R   A   F   I   A   Ș   I   S   E
  Cheia K           C   H   I   Ș   I   N   Ă   U   C   H   I   Ș   I   N   Ă   U
  Cod M             4  20  10  18  23  17   8  20   0   7  10   0  22  10  21   6
  Cod K             4   9  10  22  10  16   1  25   4   9  10  22  10  16   1  25
  M+K (mod 31)      8  29  20   9   2   2   9  14   4  16  20  22   1  26  22   0
  Criptograma C     G   Y   R   H   Â   Â   H   L   C   N   R   Ș   Ă   V   Ș   A
  ...
    C = GYRHÂÂHLCNRȘĂVȘAGBZĂÂNȚACQVYWDNUXQWFIZÂGCBÂE
>>> Decriptare cu aceeași cheie:
    M = CRIPTOGRAFIAȘISECURITATEAINFORMAȚIONALĂLAUTM   (corect: True)
\end{termout}
\end{terminal}

The final step — re-inserting spaces according to the language and the logic of the message —
is done manually: \emph{CRIPTOGRAFIA ȘI SECURITATEA INFORMAȚIONALĂ LA UTM}.

\subsection{Interactive Session and Input Validation}
\begin{terminal}
\begin{termout}
$ python3 vigenere_ro.py
======================================================================
  LAB 3 — Cifrul Vigenère pentru limba română (31 de litere)
======================================================================
  1) Criptare   2) Decriptare   0) Ieșire
  Opțiunea: 1
  Cheia (minim 7 litere): abc
[EROARE] Cheia are 3 litere; lungimea minimă este 7.
  Cheia (minim 7 litere): CHEIESECRETA
  Mesajul: Atac în zori!
[EROARE] Mesajul conține caractere nepermise: !
  Diapazonul corect: literele 'A'–'Z', 'a'–'z' și diacriticele românești
  Ă Â Î Ș Ț (31 de litere); spațiile sunt eliminate automat.
  Mesajul: Atac în zori
  Text pregătit : ATACÎNZORI
  Cheia         : CHEIESECRETA
  ...
  Criptograma: CĂELOEDSHN

  1) Criptare   2) Decriptare   0) Ieșire
  Opțiunea: 2
  Cheia (minim 7 litere): CHEIESECRETA
  Criptograma: CĂELOEDSHN
  Mesajul decriptat: ATACÎNZORI
\end{termout}
\end{terminal}

\section{Source Code and Repository}
The complete source code and this report are available in the Git repository:
\begin{center}
\url{https://github.com/don09dandidi/CS} \quad (directory \texttt{lab3/})
\end{center}

\subsection{Build and Execution Instructions}
\begin{lstlisting}[style=sh]
# Clone the repository
git clone https://github.com/don09dandidi/CS.git
cd CS/lab3

# Interactive mode (encryption / decryption)
python3 vigenere_ro.py

# Deterministic examples used in this report
python3 vigenere_ro.py --demo
\end{lstlisting}

\section{Conclusion}
\begin{enumerate}
  \item \textbf{Polyalphabetic substitution:} the Vigenère cipher was implemented for the
        31-letter Romanian alphabet strictly following
        $c_i=(m_i+k_i)\bmod 31$ and $m_i=(c_i-k_i)\bmod 31$; decryption with the same key always
        recovered the original message.
  \item \textbf{Flattened statistics:} the same plaintext letter is mapped to different
        ciphertext letters depending on its position, so simple frequency analysis (which broke
        the cipher of Lab~2) is no longer directly applicable.
  \item \textbf{Key requirements:} the minimum key length of 7 letters gives
        $31^7\approx2.75\times10^{10}$ keys; nevertheless the periodic key remains the main
        weakness — once $L$ is found by Kasiski examination or IC analysis, the cipher reduces to
        $L$ Caesar ciphers.
  \item \textbf{Robust input handling:} spaces are removed, the text is upper-cased, cedilla
        diacritics are unified, and any symbol outside the alphabet is rejected with a message
        describing the correct range.
\end{enumerate}

\begin{thebibliography}{9}
\bibitem{kahn} David Kahn, \emph{The Codebreakers: The Story of Secret Writing}, Scribner, 1996.
\bibitem{stallings} William Stallings, \emph{Cryptography and Network Security: Principles and
  Practice}, 7th Edition, Pearson, 2017.
\bibitem{guide} \emph{Lucrare de laborator nr.~3. Cifruri polialfabetice}, UTM, FCIM, 2026.
\end{thebibliography}
\end{document}
